% Part: normal-modal-logic
% Chapter: tableaux
% Section: rules-for-K

\documentclass[../../../include/open-logic-section]{subfiles}

\begin{document}

\olfileid[fr]{nml}{tab}{rul}

\olsection{Règles pour \Ax{K}}

Les règles pour les connecteurs propositionnels ordinaires sont les
mêmes que dans les !!{tableau}s propositionnels signés, auxquels on
ajoute des préfixes. Dans chaque cas, la règle appliquée à une
!!{formula} signée $\sFmla{S}{!A}[\sigma]$ produit de nouvelles
!!{formula}s portant elles aussi le préfixe~$\sigma$. Cela se
comprend intuitivement : si, par exemple, $!A \land !B$ est vraie au
monde désigné par~$\sigma$, alors $!A$ et $!B$ sont vraies en
$\sigma$ (sans que cela impose leur vérité aux autres mondes). Les
règles propositionnelles sont rassemblées dans le
\olref{tab:prop-rules}.
\iftag{prvIff}{Pour le biconditionnel primitif, $\sFmla{\True}{!A\liff!B}[\sigma]$ scinde la branche en deux : l’une reçoit $\sFmla{\True}{!A}[\sigma]$ et $\sFmla{\True}{!B}[\sigma]$, l’autre $\sFmla{\False}{!A}[\sigma]$ et $\sFmla{\False}{!B}[\sigma]$. De même, $\sFmla{\False}{!A\liff!B}[\sigma]$ produit une branche avec $\sFmla{\True}{!A}[\sigma]$ et $\sFmla{\False}{!B}[\sigma]$, et une autre avec $\sFmla{\False}{!A}[\sigma]$ et $\sFmla{\True}{!B}[\sigma]$. Chaque branche reçoit les deux formules indiquées au même préfixe.}{}
\begin{table}
  \[\def\arraystretch{3}\begin{array}{|c|c|}
    \hline
    \AxiomC{\sFmla{\True}{\lnot !A}[\sigma]}
    \RightLabel{\TRule{\True}{\lnot}}
    \UnaryInfC{\sFmla{\False}{!A}[\sigma]}
    \DisplayProof
    &
    \AxiomC{\sFmla{\False}{\lnot !A}[\sigma]}
    \RightLabel{\TRule{\False}{\lnot}}
    \UnaryInfC{\sFmla{\True}{!A}[\sigma]}
    \DisplayProof
    \\[1ex]
    \hline
    \AxiomC{\sFmla{\True}{!A \land !B}[\sigma]}
    \RightLabel{\TRule{\True}{\land}}
    \UnaryInfC{\sFmla{\True}{!A}[\sigma]}
    \noLine
    \UnaryInfC{\sFmla{\True}{!B}[\sigma]}
    \DisplayProof
    &
    \AxiomC{\sFmla{\False}{!A \land !B}[\sigma]}
    \RightLabel{\TRule{\False}{\land}}
    \UnaryInfC{$\sFmla{\False}{!A}[\sigma] \quad \mid \quad
      \sFmla{\False}{!B}[\sigma]$}
    \DisplayProof
    \\[2ex]
    \hline
    \AxiomC{\sFmla{\True}{!A \lor !B}[\sigma]}
    \RightLabel{\TRule{\True}{\lor}}
    \UnaryInfC{$\sFmla{\True}{!A}[\sigma] \quad \mid \quad
      \sFmla{\True}{!B}[\sigma]$}
    \DisplayProof
    &
    \AxiomC{\sFmla{\False}{!A \lor !B}[\sigma]}
    \RightLabel{\TRule{\False}{\lor}}
    \UnaryInfC{\sFmla{\False}{!A}[\sigma]}
    \noLine
    \UnaryInfC{\sFmla{\False}{!B}[\sigma]}
    \DisplayProof
    \\[2ex]
    \hline
    \AxiomC{\sFmla{\True}{!A \lif !B}[\sigma]}
    \RightLabel{\TRule{\True}{\lif}}
    \UnaryInfC{$\sFmla{\False}{!A}[\sigma] \quad \mid
      \quad \sFmla{\True}{!B}[\sigma]$}
    \DisplayProof
    &
    \AxiomC{\sFmla{\False}{!A \lif !B}[\sigma]}
    \RightLabel{\TRule{\False}{\lif}}
    \UnaryInfC{\sFmla{\True}{!A}[\sigma]}
    \noLine
    \UnaryInfC{\sFmla{\False}{!B}[\sigma]}
    \DisplayProof
    \\[2ex]
    \hline
  \end{array}\]
  \caption{Règles des !!{tableau}s préfixés pour les connecteurs
    propositionnels}
  \ollabel{tab:prop-rules}
\end{table}

La condition de fermeture est la même que pour les !!{tableau}s
ordinaires, à ceci près que les préfixes, et pas seulement les
!!{formula}s, doivent coïncider. Une branche est donc fermée si elle
contient à la fois
\[
\sFmla{\True}{!A}[\sigma] \quad\text{et}\quad \sFmla{\False}{!A}[\sigma]
\]
pour un certain préfixe $\sigma$ et une !!{formula}~$!A$.
\iftag{prvFalse}{Une branche contenant $\sFmla{\True}{\lfalse}[\sigma]$ est également fermée ; $\sFmla{\False}{\lfalse}[\sigma]$ n’exige aucune décomposition.}{} \iftag{prvTrue}{Une branche contenant $\sFmla{\False}{\ltrue}[\sigma]$ est également fermée ; $\sFmla{\True}{\ltrue}[\sigma]$ n’exige aucune décomposition.}{}
Les règles qui posent les hypothèses sont elles aussi celles des
!!{tableau}s ordinaires, sauf que nous attribuons toujours aux
hypothèses le préfixe~$1$. (Le choix du préfixe importe peu, pourvu
qu'il soit le même pour toutes les hypothèses.) Ainsi, nous disons que
\[
!B_1, \dots, !B_n \Proves !A
\]
si et seulement s'il existe un tableau fermé pour les hypothèses
\[
\sFmla{\True}{!B_1}[1], \dots, \sFmla{\True}{!B_n}[1],
\sFmla{\False}{!A}[1].
\]

Pour les règles modales de
\iftag{prvBox}{\iftag{prvDiamond}{$\Box$ et~}{$\Box$}}{}\iftag{prvDiamond}{$\Diamond$}{},
le préfixe de la conclusion d'une règle appliquée à !!a{formula} de
préfixe~$\sigma$ est $\sigma.n$. Toutefois, les valeurs permises de
$n$ dépendent du signe~$\True$ ou~$\False$.

\iftag{prvBox}{La règle $\TRule{\True}{\Box}$ prolonge une branche
  contenant $\sFmla{\True}{\Box !A}[\sigma]$ par
  $\sFmla{\True}{!A}[\sigma.n]$.}{}
\iftag{prvDiamond}{La règle $\TRule{\False}{\Diamond}$ prolonge une
  branche contenant $\sFmla{\False}{\Diamond !A}[\sigma]$ par
  $\sFmla{\False}{!A}[\sigma.n]$.}{}
\iftag{notprvBox,notprvDiamond}{Cette règle}{Chacune de ces règles}
ne peut s'appliquer que si le préfixe~$\sigma.n$ figure \emph{déjà}
sur la branche considérée. Nous dirons qu'un tel préfixe est
« utilisé » sur cette branche.

\iftag{prvBox}{La règle $\TRule{\False}{\Box}$ prolonge une branche
  contenant $\sFmla{\False}{\Box !A}[\sigma]$ par
  $\sFmla{\False}{!A}[\sigma.n]$.}{}
\iftag{prvDiamond}{La règle $\TRule{\True}{\Diamond}$ prolonge une
  branche contenant $\sFmla{\True}{\Diamond !A}[\sigma]$ par
  $\sFmla{\True}{!A}[\sigma.n]$.}{}
\iftag{notprvBox,notprvDiamond}{Cette règle}{Chacune de ces règles},
en revanche, ne peut s'appliquer que si le préfixe~$\sigma.n$ ne figure
\emph{pas encore} sur la branche considérée. Nous dirons qu'un tel
préfixe est « nouveau » pour cette branche.

Ces règles figurent dans le \olref{tab:rules-K}.

\begin{table}
  \begin{center}
    \def\arraystretch{3}\def\fCenter{}
    \begin{tabular}{|c|c|}
    \hline
    \iftag{prvBox}{
      \AxiomC{\sFmla{\True}{\Box !A}[\sigma]}
      \RightLabel{\TRule{\True}{\Box}}
      \UnaryInfC{\sFmla{\True}{!A}[\sigma.n]}
      \DisplayProof
      &
      \AxiomC{\sFmla{\False}{\Box !A}[\sigma]}
      \RightLabel{\TRule{\False}{\Box}}
      \UnaryInfC{\sFmla{\False}{!A}[\sigma.n]}
      \DisplayProof\\
      $\sigma.n$ est utilisé & $\sigma.n$ est nouveau
      \\[1ex]
      \hline}{}
    \iftag{prvDiamond}{
      \AxiomC{\sFmla{\True}{\Diamond !A}[\sigma]}
      \RightLabel{\TRule{\True}{\Diamond}}
      \UnaryInfC{\sFmla{\True}{!A}[\sigma.n]}
      \DisplayProof
      &
      \AxiomC{\sFmla{\False}{\Diamond !A}[\sigma]}
      \RightLabel{\TRule{\False}{\Diamond}}
      \UnaryInfC{\sFmla{\False}{!A}[\sigma.n]}
      \DisplayProof\\
      $\sigma.n$ est nouveau & $\sigma.n$ est utilisé
      \\[1ex]
      \hline}{}
    \end{tabular}
  \end{center}
  \caption{Les règles modales pour \Ax{K}.}
  \ollabel{tab:rules-K}
\end{table}

La condition imposant un préfixe déjà utilisé pour
\iftag{prvBox}{\TRule{\True}{\Box}}{\TRule{\False}{\Diamond}} est
indispensable. Sans elle, on considérerait à tort le !!{tableau}
suivant comme fermé :
\iftag{notprvBox,notprvDiamond}{%
  \iftag{prvBox}{%
  \begin{oltableau}
    [\pFmla{\True}{\Box \formula{A}}{1}, just = \TAss
      [\pFmla{\False}{\lnot\Box\lnot \formula{A}}{1}, just = \TAss
        [\pFmla{\True}{\formula{A}}{1.1}, just = {\TRule{\True}{\Box}[1]}
          [\pFmla{\True}{\Box\lnot\formula{A}}{1},
            just ={\TRule{\False}{\lnot}[2]}
            [\pFmla{\True}{\lnot\formula{A}}{1.1},
              just = {\TRule{\True}{\Box}[4]}
              [\pFmla{\False}{\formula{A}}{1.1},
                  just ={\TRule{\True}{\lnot}[5]}, close]
            ]
          ]
        ]
      ]
    ]
  \end{oltableau}
  }{
  \begin{oltableau}
    [\pFmla{\True}{\lnot\Diamond\lnot \formula{A}}{1}, just = \TAss
      [\pFmla{\False}{\Diamond\formula{A}}{1}, just = \TAss
        [\pFmla{\False}{\formula{A}}{1.1}, just = {\TRule{\False}{\Diamond}[2]}
          [\pFmla{\False}{\Diamond\lnot\formula{A}}{1},
            just ={\TRule{\True}{\lnot}[1]}
            [\pFmla{\False}{\lnot\formula{A}}{1.1},
              just = {\TRule{\False}{\Diamond}[4]}
              [\pFmla{\True}{\formula{A}}{1.1},
                  just ={\TRule{\True}{\lnot}[5]}, close]
            ]
          ]
        ]
      ]
    ]
  \end{oltableau}
}}{
  \begin{oltableau}
    [\pFmla{\True}{\Box \formula{A}}{1}, just = \TAss
      [\pFmla{\False}{\Diamond \formula{A}}{1}, just = \TAss
        [\pFmla{\True}{\formula{A}}{1.1}, just = {\TRule{\True}{\Box}[1]}
          [\pFmla{\False}{\formula{A}}{1.1},
            just = {\TRule{\False}{\Diamond}[2]}, close]
        ]
      ]
    ]
\end{oltableau}}

Or $\Box \formula{A} \Entails/ \Diamond \formula{A}$ : notre système
de démonstration serait donc incorrect. De même,
$\Diamond \formula{A} \Entails/ \Box \formula{A}$. Sans l'obligation
d'employer un préfixe nouveau pour
\iftag{prvBox}{\TRule{\False}{\Box}}{\TRule{\True}{\Diamond}},
on obtiendrait le tableau fermé suivant : \iftag{notprvBox,notprvDiamond}{%
  \iftag{prvBox}{%
  \begin{oltableau}
    [\pFmla{\True}{\lnot\Box\lnot \formula{A}}{1}, just = \TAss
      [\pFmla{\False}{\Box \formula{A}}{1}, just = \TAss
        [\pFmla{\False}{\formula{A}}{1.1}, just = {\TRule{\False}{\Box}[2]}
          [\pFmla{\False}{\Box\lnot\formula{A}}{1},
            just ={\TRule{\True}{\lnot}[1]}
            [\pFmla{\False}{\lnot\formula{A}}{1.1},
              just = {\TRule{\False}{\Box}[4]}
              [\pFmla{\True}{\formula{A}}{1.1},
                  just ={\TRule{\False}{\lnot}[5]}, close]
            ]
          ]
        ]
      ]
    ]
  \end{oltableau}
  }{
  \begin{oltableau}
    [\pFmla{\True}{\Diamond \formula{A}}{1}, just = \TAss
      [\pFmla{\False}{\lnot\Diamond\lnot\formula{A}}{1}, just = \TAss
        [\pFmla{\True}{\formula{A}}{1.1}, just = {\TRule{\True}{\Diamond}[1]}
          [\pFmla{\True}{\Diamond\lnot\formula{A}}{1},
            just ={\TRule{\False}{\lnot}[2]}
            [\pFmla{\True}{\lnot\formula{A}}{1.1},
              just = {\TRule{\True}{\Diamond}[4]}
              [\pFmla{\False}{\formula{A}}{1.1},
                  just ={\TRule{\True}{\lnot}[5]}, close]
            ]
          ]
        ]
      ]
    ]
  \end{oltableau}
}}{
  \begin{oltableau}
    [\pFmla{\True}{\Diamond \formula{A}}{1}, just = \TAss,name=one
      [\pFmla{\False}{\Box \formula{A}}{1}, just = \TAss, name=two
        [\pFmla{\True}{\formula{A}}{1.1}, just = {\TRule{\True}{\Diamond}[1]}
          [\pFmla{\False}{\formula{A}}{1.1},
            just = {\TRule{\False}{\Box}[2]}, close]
        ]
      ]
    ]
  \end{oltableau}}
\begin{explain}
Les règles propositionnelles décomposent les formules au monde
désigné par le préfixe. Les règles modales qui exigent un préfixe
successeur déjà utilisé exploitent une relation d'accessibilité
déjà établie. Celles qui exigent un préfixe nouveau introduisent
un témoin dont on peut choisir le monde.
\end{explain}
\end{document}
